\section{Major Changes (MC)}

\vspace{0.2cm}
Below we summarize the major changes for ``{\color{teal}Spectral Feature Augmentation for Graph Contrastive Learning and Beyond}'' compared to our NeurIPS'22 submission:

\vspace{0.2cm}
\begin{enumerate}
    \item To address the reviewer concerns about why $k=1$ is chosen for SFA we have added explanations in the  {\color{blue} Choice of Number of Iteration (k)} paragraph of  the~\nameref{sec:FA} section. We additionally highlighted that for $k=1$, based on Figure \ref{fig:sim2} our model enjoys the best rebalancing profile.
    
    \item To address reviewer's concern regarding the runtime of the proposed method, we ahve added  the~\nameref{sec:timing} section to illustrate that SFA incurs a negligible computation cost (Tables~\ref{tab:runningtime} and ~\ref{fig:runningtime}) and is significantly faster than SVD-based methods  (Table~\ref{tab:time_comp_svd}).
    
    \item To provide more contemporary models into the analysis as baselines, we have added recently published models (\ie, SUGLR, MERIT and COSTA) into Table~\ref{tab:MainResultNodeClassifcation}. 
    We additionally showcase that SFA is complementary to these methods by adding the SFA layer to them, and presenting results in  Table \ref{tab:node_classification_acc}  in the appendix.
    
    \item To show that our method is effective on the image datasets (CIFAR-10, CIFAR-100 and ImageNet-100), we have added more baselines(\ie, Siamese and SwAV) into Table~\ref{tab:image_result}.
    
    
    \item To show that SFA is effective in larger graph datasets, we have provided now empirical result on the PubMed in Table~\ref{tab:MainResultNodeClassifcation} and ogb-arxiv in Table~\ref{tab:obgarxiv}.
    
    
    \item To show SFA is robust to noisy data samples, we have added explanations and empirical results in the  {\color{blue} Robustness to Noisy Features} paragraph of the section \nameref{sec:spectrum}.
    
    
    \item To consider other balancing and spectrum perturbing (noise injecting strategies) methods other than SFA, we have now added a section comparing SFA with other spectrum feature augmentation methods (\ie, MaxExp(F), Grassman and Matrix preconception (MP)) in the {\color{blue} Other Variants of Spectral Augmentation} paragraph  of the~\nameref{sec:spectrum} section. Specifically, we investigate (i) in which part of spectrum the noise should be injected (leading or non-leading singular values) and (ii) what balancing profile (soft or binarized)  is best. 
    %
    We also present the formulation of MaxExp, Grassman and MP in the supplementary material in the ~\nameref{sec:maxexp}),~\nameref{sec:grass} and ~\nameref{sec:mp} sections.
    
    
    \item As one reviewer could not understand why our incomplete power iteration achieves desired rebalancing and noise injection, we have highlighted in Eq. \eqref{equ:FA_k=1} that the model achieves rebalancing of spectrum by subtracting a low-rank matrix from itself, \ie, $\widetilde{\mathbf{H}} \!=\! \mathbf{H}\! -\!\mathbf{H}_{\text{LowRank}}$. In the \nameref{sec:why_it_works} section, we show how for a small number $k\ll\infty$, power iteration essentially produces a low-rank matrix (instead of rank-1 matrix) following the power law rules. Most importantly, we highlight throughout the manuscript that Prop. \ref{prop:PowerIter1} and \ref{prop:PowerIterVar} describe the exact rebalancing profile and the spectral augmentation profile of our SFA, as illustrated in Fig. \ref{fig:sim}.
\end{enumerate}

{\color{teal}With the above changes, we have fully covered and addressed all key reviewers' concerns. %Below we summarize key parts of the NeurIPS'22 rebuttal.
}

\section{Response to Reviewers}

\vspace{0.4cm}
{\color{teal} Below we present only the gist from the original NeurIPS'22 rebuttal (kindly refer to Point 9 (`NeurIPS-22 reviews') in the CMT3 for the exact excerpt of the rebuttal.}

\vspace{0.4cm}
\subsection{Reviewer osGb (Rating: 6)}

\vspace{0.1cm} \noindent\textcolor{red}{\bf\textcolor{red}{Q1.}} {\em Will balancing the singular values help the model to distinguish negative pairs, and why?}

\noindent Paper `Understanding contrastive representation learning through alignment and uniformity on the hypersphere' by T. Wang and P. Isola points that \textbf{pushing away negative pairs makes the learned embeddings distributed uniformly} in the feature space (maximization of the information entropy). For example if we only push all positive samples closer, numerous embeddings converge to the single point. Given  $\mathbf{x}$ (row vector) in the feature matrix $\mathbf{X} \in \mathbb{R}^{n\times d}$ are uniformly distributed in the hidden space, $\mathbf{X}$ spans the entire space $\mathbb{R}^d$ rather than lies on some subspace $\mathbb{R}^k$ with fewer dimensions ($k\ll d$). Thus, promoting negative pairs to be distinct and encouraging orthogonal basis of $\mathbf{X}$ are two sides of the same coin, so to speak. If balancing has been achieved, this means model has already reached a uniform distribution in all directions, embeddings are well separated and the loss function has nothing to push away in any direction: whether the direction step is encouraged to be balanced or not, the gradient vanishes. We also stress that while singular values are balanced, the optimizer has freedom to choose the angle of update or embeddings . So the uniform distribution we achieve is in the sense of measuring angles between embeddings. 

\vspace{0.1cm} \noindent\textcolor{red}{\bf Q2.} {\em Is there a trade-off between performance and $k$? Will the method fail when the setting of is extremely large?} \noindent We address this question in {\color{blue}Point I} of {\color{blue}Major Changes}. 

The best trade-off between performance and $k$ is almost always achieved for $k=1$ as Prop.~\ref{prop:PowerIter1}, \ref{prop:PowerIterVar} and Fig. \ref{fig:sim} show that the balancing effect deteriorates when $k \geq 1$.

\vspace{0.1cm} \noindent\textcolor{red}{\textcolor{red}{\bf Q3.}} {\em I care about time complexity. In my understanding, the SVD operation should be processed on H for each iteration, which is very time-consuming.}
We address this question in {\color{blue}Point II} of {\color{blue}Major Changes}.

There is some misunderstanding, we do not use SVD in SFA (see Alg.~\ref{algo:1}) and we add experiment in \nameref{sec:timing} section which shows that SFA is significantly faster that SVD-based methods.

\vspace{0.4cm}
\subsection{Reviewer udkM (Rating: 7)}
\vspace{0.1cm} \noindent\textcolor{red}{\bf Q1. } {\em Are the random points kept fixed during all training iterations or resampled in each iteration? }

\noindent The random points are not fixed, as our Algorithm~\ref{algo:1} is a randomized algorithm. We re-sample the ${\bf r}^{(0)}$ per each run of Algorithm~\ref{algo:1}. That is how we impose the small variance on the balanced part of spectrum. Kindly see our Prop.~\ref{prop:PowerIter1} which analyses the behavior of Alg.~\ref{algo:1} and Eq.~\eqref{equ:FeatureAug} from the perspective of expected value since ${\bf r}^{(0)}$ is drawn from the Normal distribution. Prop ~\ref{prop:PowerIterVar} provides analysis of the variance. 

\vspace{0.1cm} \noindent\textcolor{red}{\bf Q2.} {\em How to obtain the results when no augmentation methods are used (the first line) in Table~\ref{tab:AblationStudy}?}

\noindent No augmentation methods means we adopt the self-contrast (two views are identical) on the original feature matrix. Of course, for that reason the results drop in that specific case, and SFA is not used for that first line. Kindly note that the negative sampling term is still active in such a case. Thus, the embedding space can still distribute features uniformly in the embedding space akin to analysis of `Understanding contrastive representation learning through alignment and uniformity on the hyper-sphere' by T. Wang and P. Isola. What is lost in such a self-view model is the invariance to the specific family of augmentations that normally are used to generate both views.

\vspace{0.4cm}
\subsection{Reviewer 55ZH (Rating: 6)}
\vspace{0.1cm} \noindent\textcolor{red}{\bf Q1.} {\em The experiments and baseline methods on image datasets are not enough. It is recommended to increase the experiments on images.}  
We address this question in {\color{blue}Point IV} of {\color{blue}Major Changes}. 

As the reviewer requested, we add more experiments on image datasets such as CIFAR-10, CIFAR-100 and ImageNet-100, and add baseliens such as Siamese and SwAV.

\vspace{0.1cm} \noindent\textcolor{red}{\bf Q2.} {\em Why is the performance better when takes $k$ a lower value? Can you provide a more detailed and in-depth explanation? }
We address this question in {\color{blue}Point I} of {\color{blue}Major Changes}. The best trade-off between performance and $k$ is almost always achieved for $k=1$ as Prop.~\ref{prop:PowerIter1}, \ref{prop:PowerIterVar} and Fig. \ref{fig:sim} show that the balancing effect deteriorates when $k \geq 1$. When $k$ is large, $\mathbf{H}_{\text{LowRank}}$ of our $\widetilde{\mathbf{H}}=\mathbf{H}-\mathbf{H}_{\text{LowRank}}$ simply converges to rank-1 matrix (based on the leading singular value).

\vspace{0.1cm} \noindent\textcolor{red}{\bf Q3. } {\em Many graph data have only simple structural relationships without providing valid features, or many features by themselves do not guarantee their authenticity and reliability. What adjustments need to be made to the model to ensure its feasibility? }
We address this question in {\color{blue}Point VI} of {\color{blue}Major Changes}. As the reviewer requested, we have added an experiment on noisy samples. It shows that a model equipped with SFA is more robust to the data noise. We explained that this is achieved as our rebalancing push-forward function focuses on leading singular values where the signal resides (we avoid boosting non-leading singular values where the noise  typically may reside).

\vspace{0.4cm}
\subsection{Reviewer fGxg (Rating: 2)}

Below we provide details of our NeurIPS'22 rebuttal for Reviewer fGxg. Firstly, we would like to highlight that the above three reviewers were satisfied with our NeurIPS'22 rebuttal (hence scoring us 6,7,6). The reviewer below however, in last moments of the discussion period, deleted parts of his comments and our responses to them  in OpenReview. Thus, explanations to his key question why our push-forward function takes form $\sigma_i(1-\lambda_i)$ were deleted. We did notify the AC and PCs about that strange behavior but we have never got any response back to acknowledge this situation. Nonetheless, {\color{blue}{Point VIII}} of {\color{blue}{Major Changes}} explains how we addressed this in the revised manuscript.


\vspace{0.4cm} \noindent\textcolor{red}{\bf Q1. } {\em Novelty of this work is very limited. By a closer look at the proposed incomplete power iteration, it is very similar to the existing preconditioning methods which use power iteration methods.}
%
We address this question in {\color{blue}Points I, VII and VIII} of {\color{blue}Major Changes}. 

We believe there exist some misunderstanding of what our Eq.~\eqref{equ:FeatureAug} does. Firstly, \textbf{we typically use $k=1$ iterations of Alg.~\ref{algo:1}}. Then Eq.~\eqref{equ:FeatureAug} becomes Eq.~\eqref{equ:FA_k=1}:

$$ \tilde{\mathbf{H}}=\mathbf{H}\left(\mathbf{I} - \frac{\mathbf{H}^\top\mathbf{H}\mathbf{r}^{(0)}\mathbf{r}^{(0)\top}\mathbf{H}^\top\mathbf{H}}{\|\mathbf{H}^\top\mathbf{H}\mathbf{r}^{(0)}\|_2^2}\right),$$

\noindent 
where $\mathbf{r}^{(0)}\sim\mathcal{N}(0;{\bf I})$ is critical as it provides spectral noise injection. The above expression shows there is never a need to use converged power iteration. On the contrary, we analyze/use properties of unconverged power iteration by multiplying feature matrix ${\bf H}$ with the complement of incomplete power iteration  to achieve flattening of spectrum according to Prop.~\ref{prop:PowerIter1} \textbf{with injected spectral noise} according to Prop.~\ref{prop:PowerIterVar} . Our method does not change the basis, \ie, in expectation, $\tilde{\bf H}$ and ${\bf H}$ enjoy the same basis which makes this method applicable to run over minibatches.

The major difference between the proposed so-called incomplete power iteration (I-PI) and power iteration in preconditioning techniques is that \textbf{our I-PI works only in the non-convergence regime}. Then the feature matrix ${\bf H}$ is multiplied with the complement of recovered from I-PI spectrum to achieve flattening only for the largest singular values rather than the entire spectrum. Note that \textbf{flattening the entire spectrum is undesired} as the smallest singular values most likely correspond to the noise rather than signal.

Because of the non-convergence characteristic, the initial stochastic point $\mathbf{r}(0)$ leads to two major theoretical properties:
\begin{itemize}
    \item I-PI will \textbf{inject random noise directly on the flattened part of the spectrum without alerting the orthogonal basis}. This helps the model focus on aligning the angle of the two graph views, thus achieving better alignment for GCL.
    
    \item The \textbf{complete power iteration (C-PI) simply recovers the largest singular value which is of no interest to us}. Figure 4 (b) shows that as C-PI converges, the injected noise diminishes to zero (e.g., the black curve for $k=8$). Thus, C-PI is of no use to us (it cannot perform spectral augmentation). \textbf{The I-PI achieves the perturbation effect that is most prominent within the flattened region of spectrum, thus encouraging not only spectrum balancing but also spectrum augmentation within balanced part of the spectrum}.
\end{itemize}

As detailed in previous responses, not only our Spectral Feature Augmentation by incomplete power iteration differs from the existing preconditioning methods but also \textbf{we have not found any method which would provide detailed characteristics of (i) how spectrum is balanced exactly (e.g., can these methods flatten only several leading singular values?) and (ii) how spectrum augmentation is achieved (if at all)}.

In contrast, \textbf{our Prop.~\ref{prop:PowerIter1} and~\ref{prop:PowerIterVar} provide exact analytical proofs for these properties and their quantities}.

\vspace{0.1cm}
\noindent\vspace{0.1cm} \noindent\textcolor{red}{\textcolor{red}{\bf Q2. }} {\em Authors are ignorant of matrix preconditioning techniques which reduce the condition number and thus achieve a balanced spectrum. }
We address this question in {\color{blue}Point VII} of {\color{blue}Major Changes}. 

We discuss and compare matrix preconditioning in Table~\ref{tab:spec_aug} and the \nameref{sec:mp} section. We have also noted both in the rebuttal and revised text that as MP reduce the spectral gap, they also magnify non-leading singular values which is the reason they are not good rebalancing methods. They also have no obvious ability to inject the spectrum into flattened leading singular values. Even randomized SVD in Table \ref{tab:spec_aug} is not performing better than the Grassman without noise based on ordinary SVD. As our main rebuttal illustrates, we tried to guess what MP design the reviewer might have had in mind, and we voiced this in our discussions but the reviewer was unable to provide any specific details as to what specific MP design they might have had in mind.

\vspace{0.1cm}
\noindent\textcolor{red}{\bf Q3. } {\em It is not clear how to tune the critical parameters $k$, such as for the incomplete power iteration. It is claimed that or for Table 6, but how to set k for other data sets and tasks?}
%
We address this question in {\color{blue}Point VI} of {\color{blue}Major Changes}. There is no need to tune the critical parameters $k$. Fig. \ref{fig:sim} illustrates Prop.~\ref{prop:PowerIter1} and Prop.~\ref{prop:PowerIterVar} to show the balancing effect deteriorates when $k \geq 1$. We simply use $k=1$ and in a very rare case, ablations in Table \ref{tab:EffectofK} showed that $k=2$ may be used (CiteSeer). We have explained this was simply cross-validated on the validation split.

\vspace{0.1cm}
\noindent\textcolor{red}{\bf Q4. } {\em Experimental results are very limited.}
%
We address this question in {\color{blue}Points III, V and VI} of {\color{blue}Major Changes}. We have added numerous experiments in the revision: 
\begin{itemize}
    \item We have added experiment on the large scale obg-arxiv in (Table~\ref{tab:obgarxiv}).
    \item We have shown that combining SFA with existing SOTA models (\ie, SUGRL) can further improve the accuracy in node clustering (\eg, SFA+SUGLR(74.4) \vs. SUGRL(73.2) in Cora) in Table~\ref{tab:node_classification_acc}.
    \item We have added more experiments on image classification (Table~\ref{tab:image_result}).
    \item We have also shown experimental results for other balancing push-forward functions in Table~\ref{tab:spec_aug} and illustrated these functions in Fig. \ref{fig:new_var}.
    \item We have also illustrated results on Matrix Preconditioning in Table \ref{tab:mp}.
\end{itemize} 




